\lectureheading{Lecture 13 (22/09/26)}

% --- page 01 ---

\section{Difficulty of the discrete logarithm}

How difficult is it to solve the discrete logarithm problem?
Recall from Lecture~12 that, for a prime \(p\), a primitive root \(g\) modulo \(p\), and \(h\in\units{p}\), the DLP asks for the integer \(x\) with \(1\le x\le p-1\) and \(g^{x}\equiv h\pmod{p}\), that is \(x=\ind_{g}(h)\). In a general group the same problem is the equation \(g^{x}=h\).

\begin{example}[An additive group]
Recall from Lecture~12 that solving the discrete logarithm can be easy in certain groups.
In the additive group \((\Z/p\Z,+)\), one is given \(g,h\in\Z/p\Z\) with
\(\gcdop(g,p)=1\), and one must find \(x\) such that
\[
    xg \equiv h \pmod{p}.
\]
The solution is
\[
    x \equiv g^{-1}h \pmod{p},
\]
and computing the inverse is easy. The same symbols in the multiplicative group \(\units{p}\) hide an apparently hard search.
\end{example}

% The date (22/09/26) is written on the page at this point.

How long does it take to solve the discrete logarithm problem?

\begin{remark}
Let \(G\) be a group and let \(g\in G\) be an element of order \(N\), so
\(g^N=e\). The equation \(g^x=h\) can be solved in \(O(N)\) steps, by
testing the powers of \(g\) one at a time. This bound is trivial. It is
also exponential in the input length: \(O(N)\) is of size \(e^{\log N}\).
% p.01: the marginal note beside ``Trivial'' is \(e^{\log N}\), i.e. the \(O(N)\) search is exponential in the input length.
\end{remark}

\begin{remark}
Depending on the group \(G\), one can do better. As in the additive example of Lecture~12, the congruence \(xg\equiv h\pmod{n}\) in \((\Z/n\Z,+)\) is ordinary linear arithmetic whenever \(\gcdop(g,n)=1\).
\end{remark}

% ============================================================
\section{Shanks' baby-step giant-step algorithm}
% ============================================================

% FIX p.01: the heading on the page is ``Shank's''.

\begin{keybox}
\begin{theorem}[Shanks]
Let \(G\) be a group, let \(g\in G\) be an element of order \(N\) with
\(N\ge 2\), and let \(h\in\langle g\rangle\). The discrete logarithm
\(g^x=h\) can be solved in \(O(\sqrt{N}\log N)\) steps. In particular
the running time is \(O(N^{1/2+\varepsilon})\) for every \(\varepsilon>0\).
% FIX p.01: under \(O(\sqrt{N}\log N)\) the page first writes \(O(N^{\varepsilon})\) for any \(\varepsilon>0\), strikes it out, and replaces it by \(O(N^{1/2+\varepsilon})\). The struck bound is too small; \(O(\sqrt{N}\log N)=O(N^{1/2+\varepsilon})\) is the right comparison.
\end{theorem}
\end{keybox}

% --- page 02 ---

The algorithm is as follows. Write
\[
    n = \bigl\lfloor\sqrt{N}\bigr\rfloor + 1,
\]
so that \(n>\sqrt{N}\). Form the two lists
\begin{align*}
    \text{List 1}&\colon
        e,\; g,\; g^2,\; \ldots,\; g^{n-1}
        &&\text{(baby steps),}\\[0.3em]
    \text{List 2}&\colon
        h,\; hg^{-n},\; hg^{-2n},\; \ldots,\; hg^{-(n-1)n}
        &&\text{(giant steps).}
\end{align*}
With \(u=g^{-n}\), List~2 is the same as
\[
    h,\; hu,\; hu^2,\; \ldots,\; hu^{n-1}.
\]

\begin{claim}
The two lists have an element in common. If
\[
    g^i = h\cdot g^{-jn},
\]
then \(h=g^{i+jn}\), and the discrete logarithm is
\[
    x = i+jn.
\]
\end{claim}

\begin{proof}
First, every integer \(k\) with \(0\le k\le N\) may be written
\(k=nq+r\) where \(0\le r<n\). The quotient satisfies
\[
    0\le q < n,
\]
because
\[
    q = \frac{k-r}{n} \le \frac{N}{n} < n.
\]
The last inequality holds since \(n>\sqrt{N}\).

Since \(h\in\langle g\rangle\), one has \(h=g^k\) for some
\(0\le k<N\). Write \(k=nq+r\) with \(0\le q,r<n\), as above. Then
\[
    h = g^k = g^{nq+r},
\]
and therefore
\[
    h\cdot g^{-nq} = g^r.
\]
The left-hand side lies in List~2 and the right-hand side lies in
List~1, so the two lists have an element in common. In the notation of
the claim, \(i=r\) and \(j=q\).
\end{proof}

% --- page 03 ---

It remains only to count steps. List~1 is built with \(n\)
multiplications, and List~2 is built with \(n\) multiplications. Sorting
List~1 takes \(O(n\log n)\) steps. Each of the \(n\) elements of List~2
can then be tested for membership in List~1 in \(O(\log n)\) steps.
The total is
\[
    O(n\log n) = O\bigl(\sqrt{N}\log N\bigr).
\]

\begin{example}[Baby-step giant-step modulo \(31\)]
Take \(p=31\) and \(G=\Fp^\times=\units{p}\). The residue \(3\)
is a primitive root, so let \(g=3\). The order of \(g\) is
\(N=p-1=30\), and
\[
    n=\bigl\lfloor\sqrt{30}\bigr\rfloor+1=6.
\]
% FIX p.03: the page writes \(\lfloor\sqrt{31}\rfloor+1=6\). The order is \(N=30\), and \(\lfloor\sqrt{30}\rfloor+1\) is the same integer, so the lists are unchanged.
List~1, the powers \(g^0,\ldots,g^{5}\) modulo \(31\), is
\[
    1,\; 3,\; 9,\; 27,\; 19,\; 26.
\]
Here \(g^n=3^6\equiv 16\pmod{31}\), so \(g^{-n}\equiv 16^{-1}\equiv 2\pmod{31}\).

For \(h=17\), List~2 is
\[
    17,\; 17\cdot 16^{-1},\; 17\cdot 16^{-2},\; \ldots,\; 17\cdot 16^{-5},
\]
which begins
\[
    17,\; 17\cdot 2,\; 17\cdot 4,\; \ldots.
\]
The second entry is \(17\cdot 2=34\equiv 3\pmod{31}\), and \(3=g^1\) lies
in List~1. Thus
\[
    3^1 = 17\cdot 3^{-6},
\]
so \(17=3^7\) and \(x=7\).

For \(h=11\), List~2 is
\[
    11,\; 11\cdot 2,\; 11\cdot 4,\; 11\cdot 8,\; 11\cdot 16,\; 11\cdot 32,
\]
which reduces modulo \(31\) to
\[
    11,\; 22,\; 13,\; 26,\; 21,\; 11.
\]
The entry \(26\) lies in List~1, since \(3^5\equiv 26\pmod{31}\), and
\(26\equiv 11\cdot 8=11\cdot 2^3\pmod{31}\). As \(2\equiv 3^{-6}\), one has
\(2^3\equiv 3^{-18}\), and therefore
\[
    3^5 = 11\cdot 3^{-18}.
\]
Hence \(3^{23}\equiv 11\pmod{31}\), so \(x=23\).
\end{example}

% --- page 04 ---

% ============================================================
\section{The Pohlig--Hellman algorithm}
% ============================================================

% FIX p.06: the page spells the name ``Pohlig--Heuman''.

Recall from Lecture~12 that if the order of \(g\) is smooth, Pohlig--Hellman solves the discrete logarithm efficiently. Given \(g^x=h\), it is enough
to find \(x\pmod{N}\), where \(N\) is the order of \(g\).

\begin{example}[A smooth order: \(p=271\)]
Let \(p=271\), and let \(g=6\), a primitive root modulo \(271\). Then
\[
    p-1 = 270 = 2\cdot 3^3\cdot 5,
\]
so it is enough to solve the discrete logarithm modulo \(2\), modulo
\(27\), and modulo \(5\). Take \(h=145\), and write \(h=g^x\).

Since \(g\) is a primitive root, \(g^{(p-1)/2}\equiv -1\pmod{p}\), and
\[
    h^{(p-1)/2}
    = g^{x(p-1)/2}
    = (-1)^x
    =
    \begin{cases}
        1, & \text{if \(x\) is even},\\
        -1, & \text{if \(x\) is odd}.
    \end{cases}
\]
The same power is the Legendre symbol \(\bigl(\frac{h}{p}\bigr)\).
Compute
\begin{align*}
    \Bigl(\frac{h}{p}\Bigr)
    = \Bigl(\frac{145}{271}\Bigr)
    &= \Bigl(\frac{5}{271}\Bigr)\Bigl(\frac{29}{271}\Bigr)
     = \Bigl(\frac{271}{5}\Bigr)\Bigl(\frac{271}{29}\Bigr)\\
    &= \Bigl(\frac{1}{5}\Bigr)\Bigl(\frac{10}{29}\Bigr)
     = \Bigl(\frac{2}{29}\Bigr)\Bigl(\frac{5}{29}\Bigr)
     = (-1)\cdot\Bigl(\frac{4}{5}\Bigr)
     = -1.
\end{align*}
% The factor \(\bigl(\frac{1}{5}\bigr)\) is crossed out on the page once it is seen to equal \(1\).
% Filled step: \(5\equiv 29\equiv 1\pmod{4}\), so quadratic reciprocity gives the two flips; \(271\equiv 1\pmod{5}\) and \(271\equiv 10\pmod{29}\); \(10=2\cdot 5\); \(\bigl(\frac{2}{29}\bigr)=-1\) because \(29\equiv 5\pmod{8}\); and \(\bigl(\frac{5}{29}\bigr)=\bigl(\frac{29}{5}\bigr)=\bigl(\frac{4}{5}\bigr)=1\).
Thus \(x\) is odd:
\[
    x\equiv 1\pmod{2}.
\]

% --- page 05 ---

Next one finds \(x\pmod{5}\). Compare
\[
    h^{(p-1)/5}
    \qquad\text{with}\qquad
    g^{(p-1)/5},\;
    g^{2(p-1)/5},\;
    g^{3(p-1)/5},\;
    g^{4(p-1)/5},\;
    g^{5(p-1)/5}.
\]
Since \((p-1)/5=54\),
\[
    145^{54}\equiv 1\pmod{p},
\]
and therefore \(x\equiv 0\pmod{5}\).

To see the shape of this congruence, write \(x=5k+m\). Then
\(h=g^x=g^{5k+m}\), so
\[
    h^{(p-1)/5}
    = g^{(5k+m)(p-1)/5}
    = g^{m(p-1)/5},
\]
the term in \(k\) disappearing because \(g^{p-1}\equiv 1\). The value
\(1\) occurs precisely for \(m\equiv 0\pmod{5}\).

Next one calculates \(x\pmod{27}\). Expand
\[
    x = c_0 + c_1\cdot 3 + c_2\cdot 3^2,
    \qquad 0\le c_i\le 2.
\]
Modulo \(p\),
\[
    g^{(p-1)/3}\equiv 242, \qquad
    g^{2(p-1)/3}\equiv 28, \qquad
    h^{(p-1)/3}\equiv 28.
\]
The last two agree, so \(c_0=2\).

To find \(c_1\), write \(h=g^x=g^{c_0+3c_1+3^2 c_2}\). Then
\[
    \frac{h}{g^{c_0}}
    = g^{3c_1+3^2 c_2}
    = g^{3(c_1+3c_2)},
\]
% FIX p.05: the middle exponent is written \(3^2 c_1\). The next expression is \(g^{3(c_1+3c_2)}\), so the factor must be \(c_2\).
and raising to the power \((p-1)/9\) isolates the next digit:
\[
    \Bigl(\frac{h}{g^{c_0}}\Bigr)^{(p-1)/9}
    = g^{c_1(p-1)/3}.
\]
% FIX p.05: the denominator is the digit \(9\). The loop looks like a \(q\), but \(q\) is not defined here, and the identity needs \((p-1)/9\): the left side is \(g^{3(c_1+3c_2)}\), so the exponent \((p-1)/9\) produces \(g^{c_1(p-1)/3}\).
The notes say to continue from here. Carrying out that arithmetic:
\(g^{c_0}=6^2\equiv 36\), and
\[
    h\cdot g^{-c_0}\equiv 132\pmod{271},
    \qquad
    132^{(p-1)/9}=132^{30}\equiv 28
    \equiv g^{2(p-1)/3}\pmod{271},
\]
so \(c_1=2\). One further step, with \(c_0+3c_1=8\), gives
\[
    h\cdot g^{-8}\equiv 3\pmod{271},
    \qquad
    3^{(p-1)/27}=3^{10}\equiv 242
    \equiv g^{(p-1)/3}\pmod{271},
\]
so \(c_2=1\). Therefore
\[
    x\pmod{27}
    = c_0+3c_1+9c_2
    = 2+6+9
    = 17.
\]
% Filled step: the page writes ``\(\ldots\) continue'' and then the answer \(c_1=2\), \(c_2=1\).
\end{example}

% --- page 06 ---

The three congruences are
\[
    \left.
    \begin{aligned}
        x&\equiv 1 \pmod{2}\\
        x&\equiv 0 \pmod{5}\\
        x&\equiv 17\pmod{27}
    \end{aligned}
    \right\}
    \qquad
    x\equiv 125\pmod{270}.
\]
% Filled step: the page records the result of the Chinese remainder theorem and not the arithmetic. From \(x=27t+17\equiv 0\pmod{5}\) one gets \(2t\equiv 3\pmod{5}\), so \(t\equiv 4\pmod{5}\) and \(x\equiv 125\pmod{270}\). The class \(125\) is already odd, so the condition modulo \(2\) holds as well. And \(6^{125}\equiv 145\pmod{271}\).

\subsection{Primes for which the algorithm does not help}

There is a way out of the improvement above. Let \(p\) be an odd prime, so
that \(p-1\) is even, and suppose that \((p-1)/2\) itself is prime. Set
\[
    q=\frac{p-1}{2},
\]
so that \(p=2q+1\). Baby-step giant-step on the whole group, and on the
subgroup of order \(q\), then cost about \(\sqrt{p}\) and \(\sqrt{q}\)
respectively, and
\[
    \sqrt{p}=\sqrt{2q+1}\approx\sqrt{2q}=\sqrt{2}\,\sqrt{q}.
\]
The two running times are of the same size, so factoring \(p-1=2q\) does
not improve on a direct attack.

\begin{definition}[Sophie Germain prime]
A \emph{Sophie Germain prime} is a prime \(p\) such that \(2p+1\) is also
prime.
\end{definition}
In the notation above, \(q\) is a Sophie Germain prime and \(p=2q+1\) is
prime.

\subsection{The general statement}

\begin{keybox}
\begin{theorem}[Pohlig--Hellman]
Let \(G\) be a group, and suppose that whenever \(g\in G\) has order
\(p^e\) one can solve the discrete logarithm in the subgroup
\(H=\langle g\rangle\) in \(S(p^e)\) steps.
% FIX p.06: the page writes the order as \(p^2\), while the cost on the next line is \(S(p^e)\) and the sum below runs over \(p_i^{e_i}\). The statement is for a general prime power.
Let \(\hat g\in G\) be an
element of order
\[
    N = p_1^{e_1}\cdots p_t^{e_t}.
\]
Then the discrete logarithm in \(\langle\hat g\rangle\) can be solved in
\[
    \sum_{i=1}^{t} S(p_i^{e_i}) + (\log N)^{*}
\]
steps.
% p.06: the second term is \((\log N)^{*}\) on the page. Page~10 identifies it as the Chinese-remainder cost \(O(\log N)\).
\end{theorem}
\end{keybox}

By comparison, baby-step giant-step on a group of order \(N\) costs
\(\sqrt{N}\log N\) steps.

% --- page 07 ---

As a rough comparison, take \(t=2\) with \(p_1\approx p_2\) and
\(e_1=e_2\). Writing \(p\) and \(e\) for those common values, and using
baby-step giant-step as the prime-power routine, so that
\(S(p^e)\) contributes a factor \(p^{e/2}\cdot e\log p\), one gets
\begin{align*}
    S(p_1^{e_1})+S(p_2^{e_2})+(\log N)^{*}
    &= p_1^{e_1/2}\cdot e_1\log p_1
       + p_2^{e_2/2}\cdot e_2\log p_2
       + (\log N)^{*}\\
    &\approx 2\,p^{e/2}\,(\log N)^{*}.
\end{align*}
For \(N\) of size \(p^{2e}\) the direct baby-step giant-step bound
\(\sqrt{N}\log N\) is about \(p^{e}\log N\), larger by a factor about
\(p^{e/2}\).

The algorithm producing the \(y_i\) in the theorem is the following.
The input is \(\hat g^{x}=h\).
\begin{enumerate}[leftmargin=1.6em]
    \item For each \(i\), set
    \[
        g_i = \hat g^{\,N/p_i^{e_i}}
        \qquad\text{and}\qquad
        h_i = h^{\,N/p_i^{e_i}}.
    \]
    The element \(g_i\) has order \(p_i^{e_i}\). Solve the discrete
    logarithm \(g_i^{y_i}=h_i\), and let \(y_i\) be the solution.
    \item Solve the system
    \[
        \begin{cases}
            x\equiv y_1\pmod{p_1^{e_1}},\\
            \qquad\vdots\\
            x\equiv y_t\pmod{p_t^{e_t}}
        \end{cases}
    \]
    by the Chinese remainder theorem. The solution satisfies
    \(\hat g^{x}=h\).
\end{enumerate}

\lectureheading{Lecture 14 (23/09/26)}

% --- page 08 ---

Baby-step giant-step applies to an element \(g\) of order \(N\): the
discrete logarithm in \(\langle g\rangle\), namely \(g^x=h\), can be
solved in \(O(\sqrt{N}\log N)\) steps. For a typical example take
\(G=\units{p}\) with \(p-1\) even, and suppose
\[
    q=\frac{p-1}{2}
\]
is also prime. As on the previous page, Pohlig--Hellman then gives no
substantial saving.

The Pohlig--Hellman theorem is recorded again, together with the
algorithm.

\begin{theorem}
Let \(G\) be a group. Suppose that for any element \(g\) of order
\(p^e\), with \(p\) prime and \(e\ge 1\), the discrete logarithm in
\(\langle g\rangle\) can be solved in \(S(p^e)\) steps. Let \(\hat g\)
be an element whose order is \(p_1^{e_1}\cdots p_t^{e_t}\). Then the
discrete logarithm in \(\langle\hat g\rangle\) can be solved in
\[
    \sum_{i=1}^{t} S(p_i^{e_i}) + (\log N)^{*}
\]
steps, by the following algorithm. Here \(N\) is the order of \(\hat g\).
\end{theorem}

% A line struck through before step~(1) is illegible; it appears to begin ``if \(N\neq p_1\), \(p_1\neq p_2\)''.
% UNCLEAR p.08: struck-out false start immediately before the algorithm.

Let \(\hat g^{x}=h\).
\begin{enumerate}[leftmargin=1.6em]
    \item Let
    \[
        g_i=\hat g^{\,N/p_i^{e_i}},
        \qquad
        h_i=h^{\,N/p_i^{e_i}},
    \]
    and let \(y_i\) be the solution of the discrete logarithm
    \(g_i^{y}=h_i\).
    \item Solve
    \[
        \begin{cases}
            x\equiv y_1\pmod{p_1^{e_1}},\\
            \qquad\vdots\\
            x\equiv y_t\pmod{p_t^{e_t}}
        \end{cases}
    \]
    by the Chinese remainder theorem. The solution satisfies
    \(\hat g^{x}=h\).
\end{enumerate}

% --- page 09 ---

\begin{proof}
The element \(g_i\) has order \(p_i^{e_i}\). If \(\hat g^{x}=h\), then
\(g_i^{x}=h_i=g_i^{y_i}\), so the true discrete logarithm satisfies
\[
    x = y_i + p_i^{e_i}\,k
\]
for some integer \(k\), for each \(i\). Conversely, any integer \(x\) of
this form for every \(i\) is a discrete logarithm, by the calculation
that follows. Substituting the shape of \(x\),
\begin{align*}
    \bigl(\hat g^{x}\bigr)^{N/p_i^{e_i}}
    &= \hat g^{\,(y_i+p_i^{e_i}k)\,N/p_i^{e_i}}
     = \hat g^{\,y_i\,N/p_i^{e_i}}
       \cdot\bigl(\hat g^{N}\bigr)^{k}
     = \hat g^{\,y_i\,N/p_i^{e_i}}\\
    &= g_i^{y_i}
     = h_i
     = h^{\,N/p_i^{e_i}}.
\end{align*}
% p.09: an attempt to conclude \(\hat g^{x}=h\) already at this line is crossed out.
The identity
\[
    \bigl(\hat g^{x}\bigr)^{N/p_i^{e_i}} = h^{\,N/p_i^{e_i}}
\]
therefore holds for every \(i\) with \(1\le i\le t\).

The integers \(N/p_1^{e_1},\ldots,N/p_t^{e_t}\) have gcd \(1\), since
each prime \(p_i\) fails to divide \(N/p_i^{e_i}\) and these are all the
prime factors of \(N\). Hence there exist integers \(c_1,\ldots,c_t\)
such that
\[
    c_1\frac{N}{p_1^{e_1}} + \cdots + c_t\frac{N}{p_t^{e_t}} = 1.
\]
Raising the identity above to the power \(c_i\) gives
\[
    \bigl(\hat g^{x}\bigr)^{c_i N/p_i^{e_i}}
    = h^{\,c_i N/p_i^{e_i}}
\]
for each \(i\). Multiplying these \(t\) equations together, the exponent
on each side is \(\sum c_i N/p_i^{e_i}=1\), and one obtains
\(\hat g^{x}=h\).
\end{proof}

The number of steps splits in the same way as the algorithm.
\begin{itemize}[leftmargin=1.4em]
    \item Step~1, accounting for each \(i\) with \(1\le i\le t\), takes
    \[
        \sum_{i=1}^{t} S(p_i^{e_i})
    \]
    steps.
    \item Step~2 solves a system of linear congruences by the Chinese
    remainder theorem.
\end{itemize}

% --- page 10 ---

For the second step, recall the explicit solution of a Chinese remainder
system. If \(x\equiv b_i\pmod{m_i}\) for \(1\le i\le t\), with the
\(m_i\) pairwise coprime, set \(M=m_1\cdots m_t\) and \(M_i=M/m_i\). Then
\(\gcdop(m_i,M_i)=1\), so there are integers \(N_i\) and \(n_i\) with
\[
    M_i N_i + m_i n_i = 1,
\]
and in particular \(M_i N_i \equiv 1\pmod{m_i}\).
% FIX p.10: the line is \(M_i N_i+m_i n_i=1\), not \(M_i N_i+m_i=1\). The coefficient \(n_i\) is the second Bézout coefficient; only \(N_i\) is used in the sum below.
The integer
\[
    x = \sum_{i=1}^{t} b_i M_i N_i
\]
is a solution. This arithmetic takes \(O(\log N)\) steps, where \(N\) is
the product of the moduli. In the Pohlig--Hellman system one has
\(m_i=p_i^{e_i}\), \(b_i=y_i\), and \(M=N\), so step~2 requires
\(O(\log N)\) steps. This is the term written \((\log N)^{*}\) above.

The prime-power cost \(S(p_i^{e_i})\) will now be reduced to the
prime-order cost \(S(p_i)\):
\[
    S(p_i^{e_i}) = e_i\, S(p_i).
\]

\begin{theorem}[Reduction to prime order]
Let \(G\) be a group and let \(p\) be a prime. Suppose \(g\in G\)
satisfies \(\ord(g)=p\), and that for every such element the discrete
logarithm in \(\langle g\rangle\) can be solved in \(S(p)\) steps. Then
for \(\hat g\in G\) with \(\ord(\hat g)=p^e\) and some \(e\ge 1\), the
discrete logarithm \(\hat g^{x}=h\) in \(\langle\hat g\rangle\) can be
solved in \(O\bigl(e\,S(p)\bigr)\) steps.
\end{theorem}

\begingroup\samepage
Applied to the previous theorem, the discrete logarithm for an element
of order \(N=\prod_i p_i^{e_i}\) can therefore be solved in
\[
    \sum_{i} e_i\, S(p_i) + O(\log N)
\]
steps.
\endgroup

\begin{proof}
One wants to solve \(\hat g^{x}=h\). The unknown may be taken in the
range \(0\le x\le p^e-1\).
% The page writes ``expect \(0\le x<p^e\)'' and, beneath it, ``\(\le p^e-1\)''.
Expand \(x\) in base \(p\),
\[
    x = x_0 + x_1 p + x_2 p^2 + \cdots + x_{e-1}p^{e-1},
    \qquad 0\le x_i<p.
\]
% --- page 11 ---
The digits \(x_i\) are what one has to find. The element
\(\hat g^{\,p^{e-1}}\) has order \(p\).

For the first digit, 
\[
    h^{p^{e-1}}
    = \bigl(\hat g^{x}\bigr)^{p^{e-1}}
    = \hat g^{\,x_0\,p^{e-1}}
    = \bigl(\hat g^{\,p^{e-1}}\bigr)^{x_0},
\]
since every term \(x_j p^j\) with \(j\ge 1\) contributes a multiple of
\(p^e\) to the exponent. Set
\[
    h' = h^{p^{e-1}},
    \qquad
    g' = \hat g^{\,p^{e-1}},
\]
both of which are determined by the input, and \(g'\) has order \(p\).
The displayed equation is \(h'=(g')^{x_0}\), a discrete logarithm for an
element of order \(p\). It is solved in \(S(p)\) steps and returns
\(x_0\). Equivalently, the discrete logarithm in the subgroup
\(\langle\hat g^{\,p^{e-1}}\rangle\) is solved in time \(S(p)\).

For the next digit,
\begin{align*}
    h^{p^{e-2}}
    = \bigl(\hat g^{x}\bigr)^{p^{e-2}}
    &= \hat g^{\,x_0\,p^{e-2}+x_1\,p^{e-1}},\\[0.3em]
    \frac{h^{p^{e-2}}}{\hat g^{\,x_0\,p^{e-2}}}
    &= \bigl(\hat g^{\,p^{e-1}}\bigr)^{x_1}.
\end{align*}
The left-hand side is known once \(x_0\) is known, and the base on the
right has order \(p\). This discrete logarithm is solved in \(S(p)\)
steps and returns \(x_1\).

Continuing in the same way produces every digit \(x_i\). At stage \(i\)
one solves
\[
    \Bigl(
        h\,\hat g^{-(x_0+x_1 p+\cdots+x_{i-1}p^{i-1})}
    \Bigr)^{p^{e-(i+1)}}
    =
    \bigl(\hat g^{\,p^{e-1}}\bigr)^{x_i}
\]
in the subgroup of order \(p\). There are \(e\) digits, so the total
number of steps is \(e\times S(p)\). This completes the proof.
\end{proof}

% --- page 12 ---

\subsection{Comparison of the algorithms}

Baby-step giant-step applies in every case. It is an exponential-time
algorithm. Pohlig--Hellman is useful in some special cases, but it does
not improve the algorithm substantially in many cases, the safeprime
situation above being the typical example.

Index calculus is a subexponential algorithm, of shape
\(e^{\sqrt{\log N}}\), and will be discussed later.
% FIX p.12: the parenthetical is \(e^{\sqrt{\log N}}\), confirmed on the page. The usual index-calculus bound also has a \(\log\log N\) under the square root; the notes omit it. The written shape is still subexponential.

% ============================================================
\section{The RSA cryptosystem}
% ============================================================

The cryptosystem is due to R.~Rivest, A.~Shamir, and L.~Adleman.

\subsection{Factorization of integers}
The main principle is that the public modulus is a product
\[
    N=pq
\]
of two large primes.

\begin{theorem}[Bertrand's postulate]
Between \(n\) and \(2n\) there is a prime, for every integer \(n\ge 2\).
\end{theorem}

If \(p\) and \(q\) are distinct primes, then
\[
    \varphi(pq)=(p-1)(q-1).
\]
Euler's theorem gives: if \(\gcdop(a,pq)=1\), then
\[
    a^{(p-1)(q-1)}\equiv 1\pmod{pq}.
\]
Is this exponent optimal?

For the special case \(p=2\), and \(q\) an odd prime,
% FIX p.12: the page writes \(\varphi(2q)=q-1\) with no hypothesis on \(q\). This holds when \(q\) is an odd prime.
\[
    \varphi(2q)=q-1.
\]
If \(\gcdop(a,2q)=1\), then
\[
    a^{q-1}\equiv 1\pmod{2q}.
\]
This exponent is optimal.
% UNCLEAR p.12: the page ends with the parenthetical ``(primitive root'', cut off at the margin. Read as a reference to primitive roots modulo \(2q\): the group \(\units{2q}\) is cyclic of order \(q-1\), so the exponent \(q-1\) is sharp.
